If $s>1/2+k$, we can define traces of all derivatives up to order $k$. Hence there is a continuous trace operator
\begin{equation}\label{eq:higher-normal-trace-continuous-product-map}
T_k:H^s(\mathbb{R}^m)\to\prod_{j=0}^kH^{s-j-1/2}(\mathbb{R}^{m-1})
\end{equation}
such that
\begin{equation}\label{eq:higher-normal-trace-smooth-values}
T_k\phi(\bm{x}^{\prime})=\left(\phi(\bm{x}^{\prime},0),\dfrac{\partial\phi}{\partial x_m}(\bm{x}^{\prime},0),\ldots,\dfrac{\partial^k\phi}{\partial x_m^k}(\bm{x}^{\prime},0)\right)
\end{equation}
for smooth functions $\phi$. Again the inverse question of constructing a function with given trace is of interest. We have

\begin{Theorem}
\label{thm:higher-normal-trace-bounded-right-inverse}
The trace operator $T_k$ has a bounded right inverse $Z_k$.
\end{Theorem}

\begingroup
\postdisplaypenalty=10000
\renewcommand{\endProof}{\par\nobreak\hfill$\blacksquare$\par\vspace{0.1cm}}
\begin{Proof}
We first define
\begin{equation}\label{eq:higher-normal-trace-auxiliary-fourier-lifting}
\widetilde Z_l\phi:=\dfrac{x_m^l}{l!(2\pi)^{(m-1)/2}K_{s-l}}\int_{\mathbb{R}^m}e^{i\bm{\xi}\cdot\bm{x}}\hat\phi(\bm{\xi}^{\prime})\dfrac{(1+|\bm{\xi}^{\prime}|^2)^{s-l-1/2}}{(1+|\bm{\xi}|^2)^{s-l}}\,\mathrm{d}\bm{\xi},
\end{equation}
where $K_{s-l}$ is as given by \eqref{eq:trace-right-inverse-normalizing-integral}. An argument analogous to the proof of Theorem~\ref{thm:trace-bounded-right-inverse} shows that $\widetilde Z_l$ is continuous from $H^{s-l-1/2}(\mathbb{R}^{m-1})$ to $H^s(\mathbb{R}^m)$ and that, for $\phi\in\mathcal{D}(\mathbb{R}^{m-1})$,
\begin{equation}\label{eq:higher-normal-trace-auxiliary-boundary-jets}
\begin{aligned}
(\widetilde Z_l\phi)(\bm{x}^{\prime},0)&=0,\quad\dfrac{\partial\widetilde Z_l\phi}{\partial x_m}(\bm{x}^{\prime},0)=0,\quad\ldots,\\
\dfrac{\partial^l\widetilde Z_l\phi}{\partial x_m^l}(\bm{x}^{\prime},0)&=\phi(\bm{x}^{\prime}).
\end{aligned}
\end{equation}
We now construct $Z_k(\phi_0,\phi_1,\ldots,\phi_k)$ recursively by the algorithm
\begin{equation}\label{eq:higher-normal-trace-right-inverse-recursion}
\begin{gathered}
v_0=\widetilde Z_0\phi_0,\quad v_{j+1}=\widetilde Z_{j+1}\left(\phi_{j+1}-\dfrac{\partial^{j+1}v_j}{\partial x_m^{j+1}}\right)+v_j,\\
Z_k(\phi_0,\phi_1,\ldots,\phi_k)=v_k.
\end{gathered}
\end{equation}

\end{Proof}
\endgroup

We note the following corollary of the trace theorem:

\begin{Corollary}
\label{cor:fractional-order-diffeomorphism-pullback}
Let $\bm{\Phi}$ be a $k$-diffeomorphism of $\mathbb{R}^m$. Then $\bm{\Phi}^*$ is a bounded linear mapping from $H^{k-1/2}(\mathbb{R}^m)$ to itself.
\end{Corollary}

\begingroup
\postdisplaypenalty=10000
\renewcommand{\endProof}{\par\nobreak\hfill$\blacksquare$\par\vspace{0.1cm}}
\begin{Proof}
We simply extend $\bm{\Phi}$ to $\mathbb{R}^{m+1}$ by defining\[\bm{\Psi}(\bm{x}^{\prime},x_{m+1})=(\bm{\Phi}(\bm{x}^{\prime}),x_{m+1}).\]Then $\bm{\Psi}$ is a $k$-diffeomorphism of $\mathbb{R}^{m+1}$ and $\bm{\Psi}^*$ is continuous from $H^k(\mathbb{R}^{m+1})$ into itself. The rest follows by taking traces.
\end{Proof}
\endgroup
